% Part: lambda-calculus
% Chapter: lambda-definability
% Section: primitive-recursive-functions

\documentclass[../../../include/open-logic-section]{subfiles}

\begin{document}

\olfileid{lam}{ldf}{prf}
\olsection{الدوال العودية البدائية \usetoken{S}{lambda definable}}

الدوال العودية البدائية، كما تقدم، هي الدوال التي تُبنى من الدوال
الأساس~$\Zero$ و$\Succ$ و$\Proj{{\OLId{latin-ordinary}{n}}}{{\OLId{latin-ordinary}{i}}}$ بإعمال التركيب والعودية البدائية.

\begin{lem}
  \ollabel{lem:basic}
  الدوال الأساس في العودية البدائية، وهي~$\Zero$ و$\Succ$
  والإسقاطات~$\Proj{{\OLId{latin-ordinary}{n}}}{{\OLId{latin-ordinary}{i}}}$، كلها !!{lambda definable}.
\end{lem}

\begin{proof}
فالحدود التي تكون بها هذه الدوال !!{lambda defined} هي:
\begin{align*}
  \fn{Zero} & \ident \lambd[{\OLId{latin-ordinary}{a}}][\lambd[fx][{\OLId{latin-ordinary}{x}}]]\\
  \fn{Succ} & \ident \lambd[{\OLId{latin-ordinary}{a}}][\lambd[fx][{\OLId{latin-ordinary}{f}} ({\OLId{latin-ordinary}{a}} {\OLId{latin-ordinary}{f}} {\OLId{latin-ordinary}{x}})]] \\
  \fn{Proj}^{\OLId{latin-ordinary}{n}}_{\OLId{latin-ordinary}{i}} & \ident \lambd[{\OLId{latin-ordinary}{x}}_{\OLMathNumeral{0}}\dots {\OLId{latin-ordinary}{x}}_{{\OLId{latin-ordinary}{n}}-{\OLMathNumeral{1}}}][{\OLId{latin-ordinary}{x}}_{\OLId{latin-ordinary}{i}}]
\end{align*}
\end{proof}

\begin{lem}
  \ollabel{lem:comp} لتكن~${\OLId{latin-ordinary}{f}}$ ذات~${\OLId{latin-ordinary}{k}}$ من الوسائط، ولتكن
  ${\OLId{latin-ordinary}{g}}_{\OLMathNumeral{0}}, \dots, {\OLId{latin-ordinary}{g}}_{{\OLId{latin-ordinary}{k}}-{\OLMathNumeral{1}}}$ ذوات~${\OLId{latin-ordinary}{n}}$ من الوسائط؛ ولتكن هذه الدوال
  !!{lambda definable} بالحدود~${\OLId{latin-looped}{F}}$ و${\OLId{latin-looped}{G}}_{\OLMathNumeral{0}}$، \dots،~${\OLId{latin-looped}{G}}_{{\OLId{latin-ordinary}{k}}-{\OLMathNumeral{1}}}$.
  فإذا عرّفت~${\OLId{latin-ordinary}{h}}$ منها بالتركيب، كانت~${\OLId{latin-ordinary}{h}}$ !!{lambda definable}
  بالحد~${\OLId{latin-looped}{H}}$.
\end{lem}

\begin{proof}
  نأخذ للحد الذي !!{lambda define}~${\OLId{latin-ordinary}{h}}$ ما يأتي:
  \[
  {\OLId{latin-looped}{H}} \ident \lambd[{\OLId{latin-ordinary}{x}}_{\OLMathNumeral{0}} \dots {\OLId{latin-ordinary}{x}}_{{\OLId{latin-ordinary}{n}}-{\OLMathNumeral{1}}}][{\OLId{latin-looped}{F}}\, ({\OLId{latin-looped}{G}}_{\OLMathNumeral{0}} {\OLId{latin-ordinary}{x}}_{\OLMathNumeral{0}} \dots
    {\OLId{latin-ordinary}{x}}_{{\OLId{latin-ordinary}{n}}-{\OLMathNumeral{1}}}) \dots ({\OLId{latin-looped}{G}}_{{\OLId{latin-ordinary}{k}}-{\OLMathNumeral{1}}} {\OLId{latin-ordinary}{x}}_{\OLMathNumeral{0}} \dots {\OLId{latin-ordinary}{x}}_{{\OLId{latin-ordinary}{n}}-{\OLMathNumeral{1}}})]
  \]
  وصحة ذلك مما يثبت مباشرة، فنترك تفصيله تمرينًا.
\end{proof}

\begin{prob}
  أكمل إثبات \olref[lam][ldf][prf]{lem:comp} ببيان أن
  ${\OLId{latin-looped}{H}}\num{{\OLId{latin-ordinary}{n}}_{\OLMathNumeral{0}}}\dots\num{{\OLId{latin-ordinary}{n}}_{{\OLId{latin-ordinary}{n}}-{\OLMathNumeral{1}}}} \red \num{{\OLId{latin-ordinary}{h}}({\OLId{latin-ordinary}{n}}_{\OLMathNumeral{0}}, \dots,
    {\OLId{latin-ordinary}{n}}_{{\OLId{latin-ordinary}{n}}-{\OLMathNumeral{1}}})}$.
\end{prob}

وليس من شروط \olref{lem:comp} أن تكون~${\OLId{latin-ordinary}{f}}$ و${\OLId{latin-ordinary}{g}}_{\OLMathNumeral{0}}$،
\dots،~${\OLId{latin-ordinary}{g}}_{{\OLId{latin-ordinary}{k}}-{\OLMathNumeral{1}}}$ عودية بدائية؛ وإنما الشرط أن تكون كلية
و!!{lambda definable}.

\begin{lem}\ollabel{lem:prim}
  لتكن~${\OLId{latin-ordinary}{f}}$ ذات~${\OLId{latin-ordinary}{n}}$ من الوسائط، و${\OLId{latin-ordinary}{g}}$ ذات~${\OLId{latin-ordinary}{n}}+{\OLMathNumeral{2}}$ من الوسائط،
  ولتكونا !!{lambda definable} بالحدين~${\OLId{latin-looped}{F}}$ و${\OLId{latin-looped}{G}}$. فإذا كانت الدالة
  ${\OLId{latin-ordinary}{h}}$ معرّفة من~${\OLId{latin-ordinary}{f}}$ و${\OLId{latin-ordinary}{g}}$ بالعودية البدائية، كانت~${\OLId{latin-ordinary}{h}}$ أيضًا
  !!{lambda definable}.
\end{lem}

\begin{proof}
  تعريف~${\OLId{latin-ordinary}{h}}$ هو بالمعادلتين
  \begin{align*}
    {\OLId{latin-ordinary}{h}}({\OLId{latin-ordinary}{x}}_{\OLMathNumeral{1}}, \dots, {\OLId{latin-ordinary}{x}}_{\OLId{latin-ordinary}{n}}, {\OLMathNumeral{0}}) &= {\OLId{latin-ordinary}{f}}({\OLId{latin-ordinary}{x}}_{\OLMathNumeral{1}}, \dots, {\OLId{latin-ordinary}{x}}_{\OLId{latin-ordinary}{n}})\\
    {\OLId{latin-ordinary}{h}}({\OLId{latin-ordinary}{x}}_{\OLMathNumeral{1}}, \dots, {\OLId{latin-ordinary}{x}}_{\OLId{latin-ordinary}{n}}, {\OLId{latin-ordinary}{y}}+{\OLMathNumeral{1}}) & = {\OLId{latin-ordinary}{g}}({\OLId{latin-ordinary}{x}}_{\OLMathNumeral{1}}, \dots, {\OLId{latin-ordinary}{x}}_{\OLId{latin-ordinary}{n}}, {\OLId{latin-ordinary}{y}}, {\OLId{latin-ordinary}{h}}({\OLId{latin-ordinary}{x}}_{\OLMathNumeral{1}}, \dots, {\OLId{latin-ordinary}{x}}_{\OLId{latin-ordinary}{n}}, {\OLId{latin-ordinary}{y}})).
  \end{align*}
  ومعنى ذلك أن التحديث الذي تعينه الدالة~${\OLId{latin-ordinary}{g}}$ يكرر عدد~${\OLId{latin-ordinary}{y}}$ من المرات،
  على أن يبدأ من الحالة~${\OLId{latin-ordinary}{f}}({\OLId{latin-ordinary}{x}}_{\OLMathNumeral{1}}, \dots, {\OLId{latin-ordinary}{x}}_{\OLId{latin-ordinary}{n}})$. وهذا هو بعينه النمط الذي يلائم
  أعداد تشيرش، إذ حقيقتها أنها مكرّرات.

  ونقتصر، تيسيرًا للكتابة، على~${\OLId{latin-ordinary}{n}}={\OLMathNumeral{1}}$: فيكون~${\OLId{latin-ordinary}{x}}$ وسيطًا واحدًا،
  ويكون~${\OLId{latin-ordinary}{y}}$ وسيط العودية. والحد الذي تكون به~${\OLId{latin-ordinary}{h}}$ !!{lambda defined}
  هو:
  \begin{align*}
    {\OLId{latin-looped}{H}} \ident & \lambd[{\OLId{latin-ordinary}{x}}][\lambd[{\OLId{latin-ordinary}{y}}][\fn{Snd} ({\OLId{latin-ordinary}{y}}
        {\OLId{latin-looped}{D}} \tuple{\num{{\OLMathNumeral{0}}}, {\OLId{latin-looped}{F}} {\OLId{latin-ordinary}{x}}})]]
    \intertext{حيث}
    {\OLId{latin-looped}{D}} \ident & \lambd[{\OLId{latin-ordinary}{p}}][\tuple{\fn{Succ} (\fn{Fst}\, {\OLId{latin-ordinary}{p}}), ({\OLId{latin-looped}{G}} {\OLId{latin-ordinary}{x}}
          (\fn{Fst}\, {\OLId{latin-ordinary}{p}}) (\fn{Snd}\, {\OLId{latin-ordinary}{p}}))}]
  \end{align*}
  والحالة المحفوظة في أثناء التكرار زوج: أوله قيمة~${\OLId{latin-ordinary}{y}}$ في تلك
  الخطوة، وثانيه قيمة~${\OLId{latin-ordinary}{h}}$ الموافقة لها. وفي كل خطوة ننشئ زوجًا جديدًا:
  أوله خلف قيمة~${\OLId{latin-ordinary}{y}}$ الحالية، وثانيه قيمة~${\OLId{latin-ordinary}{h}}$ الجديدة التي تعينها~${\OLId{latin-ordinary}{g}}$.
  فإذا استوفينا التكرار
  أخذنا المركبة الثانية، فكانت القيمة النهائية لـ~${\OLId{latin-ordinary}{h}}$. ونحقق ذلك الآن.

  المطلوب، لكل~${\OLId{latin-ordinary}{n}}$ و${\OLId{latin-ordinary}{m}}$، أن يكون~${\OLId{latin-looped}{H}}\,\num{{\OLId{latin-ordinary}{n}}}\,\num{{\OLId{latin-ordinary}{m}}} \red
  \num{{\OLId{latin-ordinary}{h}}({\OLId{latin-ordinary}{n}},{\OLId{latin-ordinary}{m}})}$. ضع
  ${\OLId{latin-looped}{D}}_{\OLId{latin-ordinary}{n}} \ident \Subst{{\OLId{latin-looped}{D}}}{\num{{\OLId{latin-ordinary}{n}}}}{{\OLId{latin-ordinary}{x}}}$. ونثبت أولًا
  ${\OLId{latin-looped}{D}}_{\OLId{latin-ordinary}{n}}^{\OLId{latin-ordinary}{m}} \tuple{\num{{\OLMathNumeral{0}}}, {\OLId{latin-looped}{F}}\, \num{{\OLId{latin-ordinary}{n}}}} \red
  \tuple{\num{{\OLId{latin-ordinary}{m}}}, \num{{\OLId{latin-ordinary}{h}}({\OLId{latin-ordinary}{n}}, {\OLId{latin-ordinary}{m}})}}$ بالاستقراء على~${\OLId{latin-ordinary}{m}}$.

  أما عند~${\OLId{latin-ordinary}{m}}={\OLMathNumeral{0}}$ فالمطلوب أن
  ${\OLId{latin-looped}{D}}_{\OLId{latin-ordinary}{n}}^{\OLMathNumeral{0}} \tuple{\num{{\OLMathNumeral{0}}}, {\OLId{latin-looped}{F}}\, \num{{\OLId{latin-ordinary}{n}}}} \red
  \tuple{\num{{\OLMathNumeral{0}}}, \num{{\OLId{latin-ordinary}{h}}({\OLId{latin-ordinary}{n}}, {\OLMathNumeral{0}})}}$.
  لكن~${\OLId{latin-looped}{D}}_{\OLId{latin-ordinary}{n}}^{\OLMathNumeral{0}} \tuple{\num{{\OLMathNumeral{0}}}, {\OLId{latin-looped}{F}}\, \num{{\OLId{latin-ordinary}{n}}}}$ هو نفس الزوج
  $\tuple{\num{{\OLMathNumeral{0}}}, {\OLId{latin-looped}{F}}\, \num{{\OLId{latin-ordinary}{n}}}}$. ثم إن~${\OLId{latin-looped}{F}}$ !!{lambda define}s~${\OLId{latin-ordinary}{f}}$،
  فيختزل هذا إلى~$\tuple{\num{{\OLMathNumeral{0}}}, \num{{\OLId{latin-ordinary}{f}}({\OLId{latin-ordinary}{n}})}}$؛
  ولأن~${\OLId{latin-ordinary}{f}}({\OLId{latin-ordinary}{n}}) = {\OLId{latin-ordinary}{h}}({\OLId{latin-ordinary}{n}}, {\OLMathNumeral{0}})$ فهذا هو~$\tuple{\num{{\OLMathNumeral{0}}}, \num{{\OLId{latin-ordinary}{h}}({\OLId{latin-ordinary}{n}},{\OLMathNumeral{0}})}}$.

  وأما خطوة الاستقراء، فافترض
  ${\OLId{latin-looped}{D}}_{\OLId{latin-ordinary}{n}}^{\OLId{latin-ordinary}{m}} \tuple{\num{{\OLMathNumeral{0}}}, {\OLId{latin-looped}{F}}\, \num{{\OLId{latin-ordinary}{n}}}} \red
  \tuple{\num{{\OLId{latin-ordinary}{m}}}, \num{{\OLId{latin-ordinary}{h}}({\OLId{latin-ordinary}{n}}, {\OLId{latin-ordinary}{m}})}}$. نريد أن يلزم منه
  ${\OLId{latin-looped}{D}}_{\OLId{latin-ordinary}{n}}^{{\OLId{latin-ordinary}{m}}+{\OLMathNumeral{1}}} \tuple{\num{{\OLMathNumeral{0}}}, {\OLId{latin-looped}{F}}\, \num{{\OLId{latin-ordinary}{n}}}} \red
  \tuple{\num{{\OLId{latin-ordinary}{m}}+{\OLMathNumeral{1}}}, \num{{\OLId{latin-ordinary}{h}}({\OLId{latin-ordinary}{n}}, {\OLId{latin-ordinary}{m}}+{\OLMathNumeral{1}})}}$، وهذا ما تبينه الاختزالات:
  \begin{align*}
    {\OLId{latin-looped}{D}}_{\OLId{latin-ordinary}{n}}^{{\OLId{latin-ordinary}{m}}+{\OLMathNumeral{1}}} \tuple{\num{{\OLMathNumeral{0}}}, {\OLId{latin-looped}{F}}\, \num{{\OLId{latin-ordinary}{n}}}}
    & \ident {\OLId{latin-looped}{D}}_{\OLId{latin-ordinary}{n}}({\OLId{latin-looped}{D}}_{\OLId{latin-ordinary}{n}}^{{\OLId{latin-ordinary}{m}}} \tuple{\num{{\OLMathNumeral{0}}}, {\OLId{latin-looped}{F}}\, \num{{\OLId{latin-ordinary}{n}}}})\\
    & \red {\OLId{latin-looped}{D}}_{\OLId{latin-ordinary}{n}}\,\tuple{\num{{\OLId{latin-ordinary}{m}}}, \num{{\OLId{latin-ordinary}{h}}({\OLId{latin-ordinary}{n}}, {\OLId{latin-ordinary}{m}})}} \text{\qquad (بفرض الاستقراء)}\\
    & \ident (\lambd[{\OLId{latin-ordinary}{p}}][
      \tuple{
        \fn{Succ} (\fn{Fst}\, {\OLId{latin-ordinary}{p}}), ({\OLId{latin-looped}{G}} \, \num{{\OLId{latin-ordinary}{n}}} (\fn{Fst}\, {\OLId{latin-ordinary}{p}}) (\fn{Snd}\, {\OLId{latin-ordinary}{p}}))
    }]) \tuple{\num{{\OLId{latin-ordinary}{m}}}, \num{{\OLId{latin-ordinary}{h}}({\OLId{latin-ordinary}{n}},{\OLId{latin-ordinary}{m}})}}\\
    & \redone
    \tuple{\fn{Succ} (\fn{Fst}\, \tuple{\num{{\OLId{latin-ordinary}{m}}}, \num{{\OLId{latin-ordinary}{h}}({\OLId{latin-ordinary}{n}},{\OLId{latin-ordinary}{m}})}}),\\
      & \qquad ({\OLId{latin-looped}{G}} \, \num{{\OLId{latin-ordinary}{n}}} (\fn{Fst}\, \tuple{\num{{\OLId{latin-ordinary}{m}}}, \num{{\OLId{latin-ordinary}{h}}({\OLId{latin-ordinary}{n}},{\OLId{latin-ordinary}{m}})}}) (\fn{Snd}\, \tuple{\num{{\OLId{latin-ordinary}{m}}}, \num{{\OLId{latin-ordinary}{h}}({\OLId{latin-ordinary}{n}},{\OLId{latin-ordinary}{m}})}}))}\\
    & \red
    \tuple{\fn{Succ}\, \num{{\OLId{latin-ordinary}{m}}}, ({\OLId{latin-looped}{G}} \, \num{{\OLId{latin-ordinary}{n}}} \,\num{{\OLId{latin-ordinary}{m}}} \, \num{{\OLId{latin-ordinary}{h}}({\OLId{latin-ordinary}{n}},{\OLId{latin-ordinary}{m}})})}\\
    & \red \tuple{\num{{\OLId{latin-ordinary}{m}}+{\OLMathNumeral{1}}}, \num{{\OLId{latin-ordinary}{g}}({\OLId{latin-ordinary}{n}}, {\OLId{latin-ordinary}{m}}, {\OLId{latin-ordinary}{h}}({\OLId{latin-ordinary}{n}}, {\OLId{latin-ordinary}{m}}))}}
  \end{align*}
  لكن~${\OLId{latin-ordinary}{g}}({\OLId{latin-ordinary}{n}}, {\OLId{latin-ordinary}{m}}, {\OLId{latin-ordinary}{h}}({\OLId{latin-ordinary}{n}}, {\OLId{latin-ordinary}{m}})) = {\OLId{latin-ordinary}{h}}({\OLId{latin-ordinary}{n}}, {\OLId{latin-ordinary}{m}}+{\OLMathNumeral{1}})$، فتتم خطوة الاستقراء.

  بقي أن نصل هذه النتيجة بالتعريف الأول. فنحسب:
  \begin{align*}
    {\OLId{latin-looped}{H}}\,\num {\OLId{latin-ordinary}{n}}\,\num {\OLId{latin-ordinary}{m}} &\ident \lambd[{\OLId{latin-ordinary}{x}}][\lambd[{\OLId{latin-ordinary}{y}}][\fn{Snd} ({\OLId{latin-ordinary}{y}}
        (\lambd[{\OLId{latin-ordinary}{p}}].\tuple{\fn{Succ} (\fn{Fst}\, {\OLId{latin-ordinary}{p}}), ({\OLId{latin-looped}{G}}\, {\OLId{latin-ordinary}{x}}\,
          (\fn{Fst}\, {\OLId{latin-ordinary}{p}})\, (\fn{Snd}\, {\OLId{latin-ordinary}{p}}))}) \tuple{\num{{\OLMathNumeral{0}}}, {\OLId{latin-looped}{F}} {\OLId{latin-ordinary}{x}}})]]\\
    & \qquad\,\num {\OLId{latin-ordinary}{n}}\, \num {\OLId{latin-ordinary}{m}}\\
    & \red \fn{Snd} (\num {\OLId{latin-ordinary}{m}} \,
    \underbrace{(\lambd[{\OLId{latin-ordinary}{p}}].\tuple{\fn{Succ} (\fn{Fst}\, {\OLId{latin-ordinary}{p}}), ({\OLId{latin-looped}{G}} \,\num {\OLId{latin-ordinary}{n}}\,
      (\fn{Fst}\, {\OLId{latin-ordinary}{p}}) (\fn{Snd}\, {\OLId{latin-ordinary}{p}}))})}_{{\OLId{latin-looped}{D}}_{\OLId{latin-ordinary}{n}}} \tuple{\num{{\OLMathNumeral{0}}}, {\OLId{latin-looped}{F}} \num {\OLId{latin-ordinary}{n}}})\\
    & \ident \fn{Snd} (\num {\OLId{latin-ordinary}{m}} \, {\OLId{latin-looped}{D}}_{\OLId{latin-ordinary}{n}}\, \tuple{\num{{\OLMathNumeral{0}}}, {\OLId{latin-looped}{F}} \num {\OLId{latin-ordinary}{n}}})\\
    & \red \fn{Snd} \,({\OLId{latin-looped}{D}}_{\OLId{latin-ordinary}{n}}^{\OLId{latin-ordinary}{m}}  \tuple{\num{{\OLMathNumeral{0}}}, {\OLId{latin-looped}{F}} \num {\OLId{latin-ordinary}{n}}})\\
    & \red \fn{Snd} \,\tuple{\num{{\OLId{latin-ordinary}{m}}}, \num{{\OLId{latin-ordinary}{h}}({\OLId{latin-ordinary}{n}}, {\OLId{latin-ordinary}{m}})}} \\
    & \red \num{{\OLId{latin-ordinary}{h}}({\OLId{latin-ordinary}{n}},{\OLId{latin-ordinary}{m}})}.
  \end{align*}
\end{proof}


\begin{prop}
  كل دالة عودية بدائية !!{lambda definable}.
\end{prop}

\begin{proof}
  فالدوال الأساس كلها !!{lambda definable} بحسب \olref{lem:basic}؛
  ثم تفيد \olref{lem:comp} و\olref{lem:prim} أن صنف الدوال
  !!{lambda definable} مغلق تحت التركيب والعودية البدائية. ومن ثم
  يشمل كل دالة عودية بدائية.
\end{proof}

\end{document}
